\documentclass[11pt]{amsart}
\usepackage{amsmath,amssymb}
\setlength{\emergencystretch}{3em}
\tolerance=600
\usepackage{tikz-cd}
\usepackage[hidelinks]{hyperref}

\newtheorem{thm}{Theorem}[section]
\newtheorem{lem}[thm]{Lemma}
\newtheorem{cor}[thm]{Corollary}
\newtheorem{prop}[thm]{Proposition}
\theoremstyle{definition}
\newtheorem{defi}[thm]{Definition}
\newtheorem{rem}[thm]{Remark}
\newtheorem{ex}[thm]{Example}

\newcommand{\Same}{\mathrm{Same}}
\newcommand{\Top}{\mathrm{Top}}
\newcommand{\Mor}{\mathrm{Mor}}
\newcommand{\Iso}{\mathrm{Iso}}
\newcommand{\Presh}{\mathrm{Presh}}
\newcommand{\Sh}{\mathrm{Sh}}
\newcommand{\Hom}{\mathrm{Hom}}
\newcommand{\im}{\mathrm{im}}
\newcommand{\sD}{\mathcal D}
\newcommand{\sE}{\mathcal E}
\newcommand{\sJ}{\mathcal J}
\newcommand{\sK}{\mathcal K}

\begin{document}

\title[Topologies on a finite category and the sameness lattice]{Grothendieck topologies on a finite category form a sublattice of the lattice of two-out-of-three classes}
\author{Ryota Shibusawa}
\thanks{Corresponding author. E-mail: \texttt{r-shibusawa@daiichi-koudai.ac.jp}}
\address{Daiichi Institute of Technology}
\email{r-shibusawa@daiichi-koudai.ac.jp}
\keywords{Grothendieck topology, subtopos, sheafification, local isomorphism,
  two-out-of-three, separated presheaf, plus construction, finite category,
  finite monoid, idempotent ideal}
\subjclass[2020]{18F10 (primary), 18F20, 18A32, 20M12 (secondary)}

\begin{abstract}
For a small category $S$ every Grothendieck topology $J$ determines the class
$W_J$ of maps of presheaves inverted by sheafification, an element of the
complete lattice $\Same(\Presh(S))$ of all classes of morphisms containing the
identities and closed under composition and two-out-of-three. The map
$J\mapsto W_J$ preserves all meets, and it was known to preserve binary joins
for finite acyclic categories and for finite categories whose two topologies
have an idempotent intersection of ideals, but not for the poset of natural
numbers. We prove that it preserves binary joins for \emph{every} finite
category: $W_{J\vee J'}=W_J\vee W_{J'}$, the join on the right being the
smallest class closed under composition and two-out-of-three containing both.
The proof has two steps. Alternating separated quotients for the two
topologies reach, in a bounded number of steps, the quotient of a presheaf by
the equivalence relation of agreement on the join-covering sieves; and on such
separated presheaves the alternating plus-construction tower is a chain of
injections into the join sheafification, which it reaches because the join
sheafification is, by an explicit retraction of the representing bimodules, a
retract of every sufficiently late stage of the tower. The second step alone does not suffice: a monoid of
order $8$ is exhibited on which the alternating plus-construction tower never
stabilizes, although the join is preserved.
\end{abstract}

\maketitle

\section{Introduction}

Let $S$ be a small category, $\Top(S)$ the lattice of Grothendieck topologies
on $S$, and for $J\in\Top(S)$ let $a_J$ be the sheafification and
$W_J=\{f:a_Jf\text{ is an isomorphism}\}$ the class of $J$-local isomorphisms
of presheaves \cite[III.7]{MLM}. A \emph{sameness} on a category $C$ is a
class of morphisms containing the identities, closed under composition and
under two-out-of-three; samenesses form a complete lattice $\Same(C)$, whose
joins are the smallest samenesses containing the union \cite{Same}. Each
$W_J$ is a sameness, $J\mapsto W_J$ is an order-embedding
$\Top(S)\hookrightarrow\Same(\Presh(S))$ preserving all meets, and always
$W_J\vee W_{J'}\subseteq W_{J\vee J'}$ with equality iff the join on the left
is saturated \cite[Thms.~3.1, 4.2]{Topjoin}. Equality holds for finite acyclic
categories and for finite categories whose topologies have an idempotent
intersection of ideals \cite[Thms.~5.1, 6.10]{Topjoin}, and fails for the
poset $\mathbb N$ \cite[Thm.~7.1]{Topjoin}; whether it holds for all finite
categories was left open there and in \cite[Question~6.8]{Integrated}. We
answer this.

\begin{thm}\label{thm:main}
For every finite category $S$ and all topologies $J,J'$ on $S$,
\[
W_J\vee W_{J'}=W_{J\vee J'}\quad\text{in }\Same(\Presh(S)).
\]
Hence $J\mapsto W_J$ is a lattice embedding $\Top(S)\hookrightarrow\Same(\Presh(S))$.
\end{thm}

\subsection*{What is known.}
The finite-sieve dictionary identifies $\Top(S)$, for finite $S$, with the
idempotent two-sided ideals of $S$ under reverse inclusion
\cite[Thm.~5.1]{Postnikov}; the join $J_D\vee J_E$ corresponds to the largest
idempotent ideal $I'$ contained in $D\cap E$. The \emph{reduction lemma}
\cite[Lemma~4.3]{Topjoin} shows that $W_D\vee W_E=W_{I'}$ as soon as the
alternating tower of plus constructions $X\to X^{+D}\to X^{+D+E}\to\cdots$
reaches an $I'$-sheaf for every presheaf $X$, and the two positive results
above were proved by showing that it does. The tower is represented by the
alternating tensor tower $T_k=D_k\otimes\cdots\otimes D_1$ of bimodules
($\Phi_kX=\Hom(T_k,X)$); in \cite[\S6.3]{Integrated} it is shown that for
every finite category $(T_k)_k$ is pro-isomorphic to $I'\otimes I'$, the
bimodule representing $a_{I'}$. The present paper does not use that result:
the only consequence needed, that $a_{I'}X$ is a retract of the late stages of
the tower, is proved here directly by an explicit retraction
(Lemma~\ref{lem:section}).

\subsection*{Related work.}
The lattice of subtoposes, or of local operators, of a presheaf topos is
classical \cite{MLM,Elephant}; for essential localizations it is described by
Kelly and Lawvere through idempotent ideals \cite{KL89}, and every
Grothendieck topology on a finite Cauchy-complete category is rigid
\cite{Marques25}, which is the background of the object-generated ideals used
below. For $\infty$-topoi, Anel, Biedermann, Finster and Joyal describe
left-exact localizations by Grothendieck topologies \cite{ABFJ2}. All of this
concerns the lattice of \emph{saturated} classes, i.e.\ of the localizations
themselves. The question here is different: after associating to each topology
its class of local isomorphisms, does the bare two-out-of-three join, before any
saturation, already coincide with the class of local isomorphisms of the joined
topology? For finite categories the answer is yes, and the proof shows which
maps outside the canonical plus-construction units are needed.

\subsection*{What is new.}
Theorem~\ref{thm:main}, and the mechanism of its proof. The plus tower alone
does not always stabilize: Section~\ref{sec:example} gives a monoid of order
$8$ on which it never does. The missing ingredient is a second, dual tower of
\emph{separated quotients}, each of which identifies the elements of a
presheaf that agree on one of the two topologies. Three facts combine:
local isomorphisms for a topology whose ideal is generated by objects are the
maps bijective at those objects (Lemma~\ref{lem:objectwise}); the alternating
separated quotients reach, in a bounded number of steps, the quotient by
agreement on the join-covering sieves (Lemma~\ref{lem:quotient}); and on a
presheaf separated for the join topology, plus constructions preserve
separatedness, so the plus tower is a chain of injections into the join
sheafification $a_{I'}X$, which it reaches because $a_{I'}X$ is a retract of
every stage beyond $4m$ (Lemmas~\ref{lem:section} and \ref{lem:plus}). The proofs are self-contained.

\section{Preliminaries}\label{sec:prelim}

Throughout, $S$ is a finite category. Presheaves are functors
$S^{\mathrm{op}}\to\mathbf{Set}$; for $a\colon x\to c$ and $v\in X(c)$ we write
$v\cdot a:=X(a)v\in X(x)$. A \emph{two-sided ideal} of $S$ is a class of
morphisms closed under composition with arbitrary morphisms on either side; it
is \emph{idempotent} if every member is a composite of two members. By the
dictionary \cite[Thm.~5.1]{Postnikov}, topologies on $S$ correspond to
idempotent ideals: $J_D$ has as covering sieves on $c$ the sieves containing
$D(-,c)=\{a\colon x\to c\text{ in }D\}$, and $J_D\le J_E$ iff $D\supseteq E$.
For the plus construction and sheafification of $J_D$ we have
$X^{+}(c)=\Hom(D(-,c),X)$ (maps of presheaves from the sieve) and
$a_DX=X^{++}$ \cite[III.5]{MLM}.

\subsection{Cauchy completion and generating objects.}
Every element $d$ of an idempotent ideal $D$ of a finite category is a
composite $d=d_N\cdots d_1$ of members for every $N$; for $N$ larger than the
number of morphisms two partial composites $d_i\cdots d_1$, $d_j\cdots d_1$
($i<j$) coincide, so $d_i\cdots d_1=w\cdot d_i\cdots d_1$ for the endomorphism
$w=d_j\cdots d_{i+1}\in D$, and a power $g=w^r$ is an idempotent of $D$ with
$d=(d_N\cdots d_{j+1})\,g\,(d_i\cdots d_1)$. Hence $D$ is the ideal generated
by the idempotent endomorphisms it contains. Passing to the Karoubi envelope
$\sK$ of $S$ (idempotents split; presheaf categories are equivalent, and
topologies, local isomorphisms and sheafifications correspond), every
idempotent ideal $D$ becomes the ideal generated by the set of objects
\[
\sD:=\{d\in\operatorname{Ob}\sK:\ \mathrm{id}_d\in D\},
\]
in the sense that $D(-,c)$ is the set of morphisms into $c$ factoring through
an object of $\sD$. We may therefore, and do, assume that $S$ is Cauchy
complete and that $D$, $E$ are the ideals generated by the object sets
$\sD=\{d:\mathrm{id}_d\in D\}$ and $\sE=\{e:\mathrm{id}_e\in E\}$; the join
$J_D\vee J_E$ is then $J_{I'}$ where $I'$ is the ideal generated by
$\sJ:=\sD\cap\sE=\{j:\mathrm{id}_j\in D\cap E\}$ (the idempotent ideal generated
by the idempotents of $D\cap E$ is the largest idempotent ideal in $D\cap E$).
Each of the descending chains $(DE)^r$ and $(ED)^r$, $r\ge1$, stabilizes at an
idempotent ideal contained in $D\cap E$ and containing $I'=I'^{2r}$, hence at
$I'$; we fix $m\ge1$ with $(DE)^m=(ED)^m=I'$, so that every product of $2m$
alternating factors $D,E$, in either order, equals $I'$.

\subsection{Samenesses and the join theorem.}
For a functor $P$ write $W_P=\{f:Pf\text{ invertible}\}$, a sameness.

\begin{thm}[join theorem, {\cite[Thm.~4.1]{Postnikov}}]\label{thm:join}
Let $U,V$ be samenesses on $C$ whose join contains the isomorphisms of $C$
(as is the case when $U$ or $V$ is of the form $W_F$), and let $P\colon C\to C$
be a functor with a natural transformation $\pi\colon\mathrm{id}\Rightarrow P$
such that $\pi_X\in U\vee V$ for all $X$ and $P$ inverts $U$ and $V$. Then
$U\vee V=W_P$.
\end{thm}
\begin{proof}
$W_P\supseteq U\vee V$ since $W_P$ is a sameness containing $U$ and $V$. If
$Pf$ is invertible, then $Pf\in U\vee V$ by hypothesis, so the naturality
square $\pi_Y f=Pf\,\pi_X$ has three sides in $U\vee V$, hence the fourth.
\end{proof}

\subsection{The plus tower and its representing bimodules.}\label{sub:tower}
An $S$--$S$-bimodule is a functor $S^{\mathrm{op}}\times S\to\mathbf{Set}$; ideals
are sub-bimodules of $\Hom_S$, and for bimodules $A,B$ the tensor
$(A\otimes B)(x,z)=\coprod_yA(y,z)\times B(x,y)/\bigl((a\cdot m,b)\sim(a,m\cdot b)\bigr)$
is a bimodule. For a bimodule $A$ and a presheaf $X$ the presheaf
$\Hom(A,X)(c):=\Hom_{\Presh(S)}(A(-,c),X)$ satisfies the adjunction
$\Hom(A\otimes B,X)\cong\Hom(A,\Hom(B,X))$ and, for $A=\Hom_S$, $\Hom(\Hom_S,X)\cong X$;
a map of bimodules $\alpha\colon A\to A'$ induces $\alpha^*\colon\Hom(A',X)\to\Hom(A,X)$,
and by the Yoneda lemma $\alpha$ is an isomorphism iff $\alpha^*$ is bijective
for all $X$. Elements of $(A_k\otimes\cdots\otimes A_1)(x,c)$ are classes
$[a_k,\dots,a_1]$ of composable chains $x\to\cdots\to c$ with $a_i\in A_i$,
modulo moving a morphism of $S$ across a tensor sign.

Write $D_1=D$, $D_2=E$, $D_3=D,\dots$. Since $D(-,c)$ is the least
$J_D$-covering sieve, $X^{+D}=\Hom(D,X)$, with unit $x\mapsto(d\mapsto x\cdot d)$
induced by the multiplication $D\to\Hom_S$. Hence the alternating plus tower is
\[
\Phi_kX:=X^{+D_1+D_2\cdots+D_k}=\Hom(T_k,X),\qquad T_k:=D_k\otimes\cdots\otimes D_1,
\]
with unit $u_k\colon X\to\Phi_kX$, $u_k(x)=(t\mapsto x\cdot\bar t)$, where
$\bar t=d_k\cdots d_1$ is the composite of the chain $t$; $u_k$ is a composite
of plus-construction units, hence of maps of $W_D\cup W_E$. Likewise
$a_{I'}X=X^{+I'+I'}=\Hom(P,X)$ for $P:=I'\otimes I'$, with unit
$\eta_X(x)=(p\mapsto x\cdot\bar p)$.

\begin{lem}\label{lem:mult}
Let $m_k\colon P\otimes T_k\to P$ be the bimodule map
$[i,i',d_k,\dots,d_1]\mapsto[i,\,i'd_k\cdots d_1]$ (well defined since $I'$ is a
right ideal). Then $m_k$ is an isomorphism of bimodules, and
$m_k^*=a_{I'}(u_k)\colon a_{I'}X\to a_{I'}\Phi_kX$ under the identification
$a_{I'}\Phi_kX=\Hom(P,\Hom(T_k,X))=\Hom(P\otimes T_k,X)$.
\end{lem}
\begin{proof}
For $\psi\in a_{I'}X=\Hom(P,X)$ the map $a_{I'}(u_k)(\psi)=u_k\circ\psi$ sends
$p$ to $(t\mapsto\psi(p)\cdot\bar t)=(t\mapsto\psi(p\cdot\bar t))$, i.e.\
corresponds to $[p,t]\mapsto\psi(m_k[p,t])$; so $a_{I'}(u_k)=m_k^*$. Since
$u_k\in W_D\vee W_E\subseteq W_{I'}$ ($J_D,J_E\le J_{I'}$), $a_{I'}(u_k)$ is
bijective for every $X$, and $m_k$ is an isomorphism by Yoneda.
\end{proof}

Let $\iota_k\colon\Phi_kX\to a_{I'}X$ be the sheafification unit of $\Phi_kX$
followed by $(a_{I'}u_k)^{-1}$; thus $\iota_k=(m_k^{*})^{-1}\circ\varepsilon_k^*$
where $\varepsilon_k\colon P\otimes T_k\to T_k$, $[p,t]\mapsto\bar p\cdot t$, is
the bimodule map inducing the unit $\Phi_kX\to a_{I'}\Phi_kX$. By naturality of
units, $\iota_ku_k=\eta_X$.

\begin{lem}[retraction]\label{lem:section}
Let $k\ge4m$ and choose $j$ with $2m\le j\le k-2m$. The map
\[
r_k\colon T_k\to P,\qquad[d_k,\dots,d_1]\mapsto[d_k\cdots d_{j+1},\ d_j\cdots d_1],
\]
is a well-defined map of bimodules (both products lie in $I'$, being products
of at least $2m$ alternating factors), and $r_k\circ\varepsilon_k=m_k$.
Consequently $\sigma_k:=r_k^*\colon a_{I'}X\to\Phi_kX$ satisfies
$\iota_k\sigma_k=\mathrm{id}$; in particular $\iota_k$ is surjective for every
presheaf $X$ and every $k\ge4m$.
\end{lem}
\begin{proof}
A tensor move inside one of the two blocks does not change its product; a move
across the boundary, $[\dots,d_{j+1}m,d_j,\dots]\sim[\dots,d_{j+1},md_j,\dots]$,
changes the pair of products from $(Am,B)$ to $(A,mB)$, which are identified in
$P=I'\otimes I'$. So $r_k$ is well defined, and it is clearly a bimodule map.
For $[p,t]\in P\otimes T_k$ with $p=[i,i']$,
$r_k\varepsilon_k[p,t]=[ii'd_k\cdots d_{j+1},\,d_j\cdots d_1]
=[i,\,i'd_k\cdots d_1]=m_k[p,t]$, the middle equality being a tensor move in
$P$ (both $i$ and $d_j\cdots d_1$ lie in $I'$). Hence
$\iota_k\sigma_k=(m_k^*)^{-1}\varepsilon_k^*r_k^*=(m_k^*)^{-1}(r_k\varepsilon_k)^*=\mathrm{id}$.
\end{proof}

\section{Local isomorphisms are objectwise bijections}\label{sec:objectwise}

\begin{lem}\label{lem:objectwise}
Let $D$ be the ideal generated by a set of objects $\sD$ of the Cauchy complete
category $S$. A map of presheaves $f\colon X\to Y$ is a $J_D$-local isomorphism
iff $f_d\colon X(d)\to Y(d)$ is bijective for every $d\in\sD$. Consequently
$W_D=\{f:f\text{ bijective at }\sD\}$, $W_E=\{f:f\text{ bijective at }\sE\}$
and $W_{I'}=\{f:f\text{ bijective at }\sJ\}$.
\end{lem}
\begin{proof}
$f$ is $J_D$-locally injective iff for all $c$ and $x,x'\in X(c)$ with
$fx=fx'$ the sieve $\{a:x\cdot a=x'\cdot a\}$ contains $D(-,c)$, and
$J_D$-locally surjective iff for all $y\in Y(c)$ the sieve
$\{a:y\cdot a\in\im f\}$ contains $D(-,c)$ \cite[III.7]{MLM}. Every
$a\in D(-,c)$ factors as $a=p\circ q$ with $p\colon d\to c$, $d\in\sD$. If
$f_d$ is injective, then $f(x\cdot p)=(fx)\cdot p=(fx')\cdot p=f(x'\cdot p)$
gives $x\cdot p=x'\cdot p$, hence $x\cdot a=x'\cdot a$; conversely local
injectivity with $c=d$ and $a=\mathrm{id}_d$ gives injectivity at $d$. If $f_d$
is surjective then $y\cdot p=f(x)$ for some $x\in X(d)$ and
$y\cdot a=f(x\cdot q)$; conversely local surjectivity with $a=\mathrm{id}_d$
gives surjectivity at $d$. The last assertion uses that $I'$ is generated by
$\sJ$.
\end{proof}

For a finite monoid $M$ and an idempotent two-sided ideal $D$: a map of right
$M$-sets $f$ is a $J_D$-local isomorphism iff $f|_{Xe}\colon Xe\to Ye$ is
bijective for every idempotent $e\in D$.

\section{Separated quotients}\label{sec:quotient}

\begin{defi}
For a presheaf $X$ and an ideal $D$ generated by objects $\sD$, let
$x\approx_Dx'$ (for $x,x'\in X(c)$) iff $x\cdot a=x'\cdot a$ for all
$a\in D(-,c)$; equivalently, iff $x\cdot p=x'\cdot p$ for all $p\colon d\to c$
with $d\in\sD$. The \emph{$D$-separated quotient} is $Q_DX:=X/{\approx_D}$.
\end{defi}

If $x\approx_Dx'$ then $x\cdot b\approx_Dx'\cdot b$ for every $b\colon c'\to c$
(as $(x\cdot b)\cdot a=x\cdot(ba)$ and $ba\in D(-,c)$), so $Q_DX$ is a presheaf
and $q\colon X\twoheadrightarrow Q_DX$ is a map; $Q_DX$ is $J_D$-separated, and
$q$ is bijective at every $d\in\sD$ (take $a=\mathrm{id}_d$), so $q\in W_D$ by
Lemma~\ref{lem:objectwise}.

\begin{lem}\label{lem:quotient}
Let $X_0=X$ and $X_{k+1}=Q_{D_{k+1}}X_k$ with the alternation $D_1=D$, $D_2=E,\dots$.
Then two elements of $X(c)$ have the same image in $X_k$ iff they agree on
restriction along every morphism of the product ideal $D_k\circ\cdots\circ D_1$;
in particular $X_k=Q_{I'}X$ for $k\ge2m$, where $Q_{I'}X$ is the quotient by
agreement along all morphisms from objects of $\sJ$. The composite
$X\twoheadrightarrow Q_{I'}X$ lies in $W_D\vee W_E$.
\end{lem}
\begin{proof}
Induction on $k$: $x,x'\in X(c)$ have the same image in $X_{k+1}$ iff
$x\cdot b$ and $x'\cdot b$ have the same image in $X_k$ for all
$b\in D_{k+1}(-,c)$, i.e.\ iff $x\cdot(ba)=x'\cdot(ba)$ for all such $b$ and
all $a\in(D_k\cdots D_1)(-,\operatorname{dom}b)$; the morphisms $ba$ are
exactly those of $D_{k+1}D_k\cdots D_1$. For $k\ge2m$ this ideal is $I'$, and
agreement along $I'$ is agreement along the morphisms $p\colon j\to c$ with
$j\in\sJ$ (every $a\in I'(-,c)$ is $p\circ q$ with such a $p$). Each quotient
map $X_k\to X_{k+1}$ lies in $W_{D_{k+1}}$.
\end{proof}

\section{The plus tower on a separated presheaf}\label{sec:plus}

Call a presheaf \emph{$\sJ$-separated} if it is $J_{I'}$-separated, i.e.\
$x\cdot p=x'\cdot p$ for all $p\colon j\to c$, $j\in\sJ$, implies $x=x'$.

\begin{lem}\label{lem:plus}
Let $Y$ be $\sJ$-separated and $\sJ\subseteq\sD$. Then $Y^{+D}$ is
$\sJ$-separated. Consequently every stage $\Phi_kY$ of the alternating plus
tower is $\sJ$-separated, the maps $\iota_k\colon\Phi_kY\to a_{I'}Y$ are
injective, and for $k\ge4m$ they are bijective; the unit $Y\to a_{I'}Y$ lies
in $W_D\vee W_E$.
\end{lem}
\begin{proof}
For $j\in\sJ\subseteq\sD$ the sieve $D(-,j)$ is maximal, so $Y^{+D}(j)=Y(j)$,
and for $\varphi\in Y^{+D}(c)=\Hom(D(-,c),Y)$ and $p\colon j\to c$ the
restriction is $\varphi\cdot p=\varphi(p)$. Let $\varphi,\psi\in Y^{+D}(c)$
satisfy $\varphi(p)=\psi(p)$ for all $p\colon j\to c$ with $j\in\sJ$. For any
$n\colon x\to c$ in $D(-,c)$ and any $b\colon j\to x$ with $j\in\sJ$,
$\varphi(n)\cdot b=\varphi(nb)=\psi(nb)=\psi(n)\cdot b$, and $Y$ is
$\sJ$-separated, so $\varphi(n)=\psi(n)$; thus $\varphi=\psi$. The same holds
for $E$. Hence all $\Phi_kY$ are $\sJ$-separated. The map $\iota_k$ is the
sheafification unit of $\Phi_kY$ followed by an isomorphism
(Section~\ref{sub:tower}), and the unit of a separated presheaf is injective
\cite[III.5]{MLM}; so $\iota_k$ is injective. By Lemma~\ref{lem:section},
$\iota_k$ is surjective for $k\ge4m$. So $\iota_k$ is an isomorphism; isomorphisms lie in
$W_D\cap W_E\subseteq W_D\vee W_E$ (a sameness need only contain the
identities, but $W_D$ and $W_E$ are classes of maps inverted by a functor), so
the unit $Y\to a_{I'}Y=\iota_ku_k$ lies in $W_D\vee W_E$.
\end{proof}

\begin{figure}[ht]
\[
\begin{tikzcd}[column sep=large]
X\arrow[r,"W_D",two heads]\arrow[rrrr,bend left=22,"\eta_X"] & Q_DX\arrow[r,"W_E",two heads] & \cdots\arrow[r,two heads] & Q_{I'}X\arrow[r,"u_k\ \in W_D\vee W_E"] & \Phi_kQ_{I'}X\arrow[d,"\iota_k","\simeq"']\\
& & & & a_{I'}X
\end{tikzcd}
\]
\caption{The proof of Theorem~\ref{thm:main}: separated quotients (at most $2m$)
followed by plus constructions ($4m$) connect every presheaf to its join
sheafification inside the join.}
\end{figure}

\section{Proof of the theorem}\label{sec:proof}

\begin{proof}[Proof of Theorem~\ref{thm:main}]
By Section~\ref{sec:prelim} we may assume $S$ Cauchy complete, $J=J_D$,
$J'=J_E$ with $D,E$ generated by objects, and $J\vee J'=J_{I'}$. Apply the
join theorem with $U=W_D$, $V=W_E$, $P=a_{I'}$ and $\pi$ the unit
$\eta\colon\mathrm{id}\Rightarrow a_{I'}$. $P$ inverts $W_D$ and $W_E$ since
$J_D,J_E\le J_{I'}$. For a presheaf $X$ the unit factors as
\[
X\twoheadrightarrow Q_{I'}X\xrightarrow{\ \eta\ }a_{I'}Q_{I'}X\cong a_{I'}X,
\]
the first map in the join by Lemma~\ref{lem:quotient}, the second by
Lemma~\ref{lem:plus} ($Q_{I'}X$ is $\sJ$-separated), and the isomorphism
induced by the quotient map, which $a_{I'}$ inverts. Hence
$\eta_X\in W_D\vee W_E$ for all $X$, and $W_D\vee W_E=W_{a_{I'}}=W_{I'}$.
Since $W_{J\vee J'}=\mathrm{Sat}(W_J\vee W_{J'})\supseteq W_J\vee W_{J'}$
always, and $J\mapsto W_J$ preserves meets \cite[Thm.~3.1]{Topjoin}, the map
is a lattice embedding.
\end{proof}

\begin{rem}
The number of steps is uniform: $2m$ quotients and $4m$ plus constructions,
depending only on $(S,D,E)$. The proof uses finiteness of $S$ only through the
existence of $m$, the stabilization of the alternating products of the two
ideals at $I'$; this fails for $\mathbb N$, where the join is not preserved
\cite[Thm.~7.1]{Topjoin}. The number $2m$ of quotients is also the number of
alternating iterations of the two Lawvere--Tierney local operators needed to
reach their join, since $j_D j_E(R)=\{a: a\cdot DE\subseteq R\}$ on sieves $R$.
\end{rem}

\begin{rem}
Nothing in the argument requires $S$ to be Cauchy complete except the
description of ideals by objects; for a finite category $S$ the statement is
transported along the equivalence $\Presh(S)\simeq\Presh(\sK)$.
\end{rem}

\section{The plus tower alone does not suffice}\label{sec:example}

The reduction lemma of \cite{Topjoin} proves join preservation whenever the
alternating plus tower reaches an $I'$-sheaf. The following example shows that
Lemma~\ref{lem:quotient} cannot be dispensed with.

\begin{prop}\label{prop:nil}
Let $\overline M=\{1,1_a,1_b,u,v,uv,vu,0\}$ be the contracted monoid of the
category with two objects $a,b$, arrows $u\colon a\to b$, $v\colon b\to a$, in
which all composites of length $\ge3$ are identified with $0$; $|\overline M|=8$.
Let $D$ be the ideal generated by $1_b$ and $E$ the ideal generated by $1_a$;
then $D\cap E=\{u,v,uv,vu,0\}$ is not idempotent and $I'=\{0\}$. The
alternating plus-construction tower of $(D,E)$ never stabilizes: for every
even $k$ the composite multiplication map $T_{k+4}\to T_k$ of the
representing tensor tower is not surjective. In particular the maps
$T_{k+1}\to T_k$ are not eventually bijective, and there is no $k$ with
$\Phi_kX$ an $I'$-sheaf for all $X$. Nevertheless $W_D\vee W_E=W_{I'}$ by
Theorem~\ref{thm:main}.
\end{prop}
\begin{proof}
Here $m=2$, so the image of the composite $T_{k+4}\to T_k$ is
$(D_{k+4}\cdots D_{k+1})\cdot T_k=I'\cdot T_k$ (Section~\ref{sec:prelim}),
the set of classes with a representative whose first slot lies in
$I'=\{0\}$. We show that $T_k\ne I'\cdot T_k$ for every even $k$; the last
two assertions then follow, since bijectivity of $\mu_{k+1}$ for all $k\ge K$
would make $T_{K+4}\to T_K$ bijective, and since, if $\Phi_kX$ were an
$I'$-sheaf for all $X$, then $\Phi_kX$ would be a $D$-sheaf and an $E$-sheaf
($J_D,J_E\le J_{I'}$), all later plus-construction units on it would be
isomorphisms and $\Phi_kX\to\Phi_{k+l}X$ would be bijective for all $l$ and
$X$: for even $k$ this makes $T_{k+4}\to T_k$ an isomorphism of bimodules by
Yoneda, and for odd $k$ it makes $\Phi_{k+1}X\cong\Phi_kX$ an $I'$-sheaf for
all $X$ with $k+1$ even. Consider, for even $k=2j$, the class
of $t=(1_a,uv,1_a,uv,\dots)\in T_k$, whose slots alternate between $E$ and $D$
as required. Tuples of non-zero elements with product the word $w=(uv)^j$ of
the free category are \emph{cut configurations} of $w$: cut points
$0=c_0\le\dots\le c_k=2j$, the $i$-th slot being the subword between $c_{i-1}$
and $c_i$, an empty slot being the identity at its cut point, which must be
$a$ (even cut point) for an $E$-slot ($i$ odd) and $b$ (odd cut point) for a
$D$-slot ($i$ even); the slots have length $\le2$, since longer words are $0$.
A tensor move between tuples of non-zero elements moves one cut point. Put
$d_i=i-c_i$; a slot of length $\ell$ changes $d$ by $1-\ell$, and an empty
slot raises $d$ by one and is allowed only from an even value of $d$ (the
parity condition on $c_i$). From an odd value $d$ cannot rise, so $d\le1$
throughout; a slot of length $\ge3$ would send $d$ to a value $\le-1$, from
which $d_k=0$ can no longer be reached. Hence no cut configuration has a slot
of length $\ge3$, and no tensor move applied to a cut configuration produces a
zero slot (a zero would have to arise as a composite of length $\ge3$ or as a
non-composable pair, and adjacent slots of a configuration are composable). So
the class of $t$ consists of cut configurations only and has no representative
with a zero slot; thus $T_k\ne I'\cdot T_k$ for every even $k$.
\end{proof}

\begin{rem}
Enumerating the tensor relations gives $|T_1|=6$ and $|T_k|=5$ for $k\ge2$:
besides the zero class, four classes of cut configurations of non-zero words
persist at every level, and no multiplication map is surjective.
\end{rem}

By contrast, the separated quotients collapse these classes at once: on
$\overline M$-sets the quotient $Q_E$ identifies two elements that agree
under $1_a$, $u$, $v$, $uv$, $vu$ and $0$, and since $m=2$ here, four
alternating quotients produce the $I'$-separated quotient, which is already the
$I'$-sheaf: the set of elements fixed by $0$. The derivation of the unit
$X\to a_{I'}X$ inside $W_D\vee W_E$ is thus a composite of quotient maps,
none of which is a plus-construction unit.

\section*{Funding}
No funding was received for this work.

\section*{Declarations}
\noindent\textbf{CRediT authorship contribution statement.}\quad Ryota Shibusawa:
Conceptualization, Methodology, Formal analysis, Investigation, Software,
Validation, Visualization, Writing -- original draft, Writing -- review \& editing.

\noindent\textbf{Competing interests.}\quad The author declares that he has no
competing interests.

\begin{thebibliography}{9}
\bibitem{ABFJ2} M.~Anel, G.~Biedermann, E.~Finster, A.~Joyal, \emph{Left-exact
localizations of $\infty$-topoi II: Grothendieck topologies}, J. Pure Appl.
Algebra \textbf{228} (2024), 107597.
\bibitem{Elephant} P.~T.~Johnstone, \emph{Sketches of an Elephant: A Topos
Theory Compendium}, Oxford University Press, 2002.
\bibitem{KL89} G.~M.~Kelly, F.~W.~Lawvere, \emph{On the complete lattice of
essential localizations}, Bull. Soc. Math. Belg. S\'er. A \textbf{41} (1989),
289--319.
\bibitem{MLM} S.~Mac Lane, I.~Moerdijk, \emph{Sheaves in Geometry and Logic},
Springer, 1992.
\bibitem{Marques25} J.~Marqu\`es, \emph{A criterion for categories on which
every Grothendieck topology is rigid}, Appl. Categ. Structures \textbf{33}
(2025), art.\ 37.
\bibitem{Same} R.~Shibusawa, \emph{The sameness lattice of a category:
weak-equivalence structures, Galois functoriality, and comparison intervals},
preprint (2026), archived at DOI 10.5281/zenodo.22823772.
\bibitem{Postnikov} R.~Shibusawa, \emph{Postnikov truncations as joins of sheaf
and homotopy equivalences}, preprint (2026), archived at DOI 10.5281/zenodo.22956181.
\bibitem{Topjoin} R.~Shibusawa, \emph{Meets and joins of Grothendieck
topologies in the sameness lattice}, preprint (2026), archived at DOI
10.5281/zenodo.22978446.
\bibitem{Integrated} R.~Shibusawa, \emph{Sheaf equivalences, Postnikov
truncations and preorder reflections in the lattice of two-out-of-three
classes}, preprint (2026), archived at DOI 10.5281/zenodo.23200684.
\end{thebibliography}

\end{document}
